\section{The Boundary: Structure and Experience (P9)}
\label{sec:mind}

Vibe Theory is, at root, a theory of experience. The vibe is a unit of experience, and the axiom is that to exist is to experience or be experienced. P9 asks what the testbed can and cannot say about this.

What it can say is structural. We can locate, on the mesh, where a self would sit on the framework's own terms: a region that is well screened from its surroundings (a Markov blanket) and highly integrated within itself.

\begin{result}[Structural correlates of a self]
On a hyperbolic mesh with random tones the testbed measures a Markov-blanket screening score of $0.77$ and an integration proxy of $1.91$, locating a well-screened, well-integrated region, the structural signature the framework predicts a self would have.
\end{result}

What it cannot say is whether those correlates \emph{are} experience. That a region is well screened and highly integrated is a fact about structure. Whether it feels like anything is the framework's foundational claim, not a measurement the simulator can make. This is the honest boundary, and it is by design.

The point is worth stating plainly, because it is where Vibe Theory differs from theories that try to derive consciousness. Vibe Theory does not derive experience. It posits experience as the foundation and derives the physical from it. So the testbed is the right tool for the physical ladder, geometry through force, and the wrong tool for the foundational claim. We measure the structure. The experiential reading is the framework's jurisdiction, and we leave it there rather than dress a structural number as a solution to the hard problem of consciousness (Chalmers, 1995).
